% Appendices. Material: docs/requirements/traceability.md, docs/research/traceability.md,
% curriculum/lexicon.ru.yaml (transcription), docs/research/preregistration.md.
\chapter{Traceability}\label{app:trace}

\section{Research traceability}

Every number in Chapter~\ref{ch:eval} is linked to a research question, a hypothesis, a GQM question, the
code that computes it and the figure or table where it appears.

\begin{table}[ht]
\centering\small
\caption{Research question $\to$ hypothesis $\to$ metric $\to$ computation $\to$ result.}
\begin{tabularx}{\linewidth}{@{}lllLl@{}}
\toprule
RQ & Hyp. & GQM & Metric (computed by) & Result \\
\midrule
RQ-1 & H1.1 & Q1.1--Q1.2 & precision, recall, F1 (metrics module) & T-01 \\
RQ-1 & H1.2 & Q1.3 & FAR, FRR, agreement with the tutor (metrics module, SQL report 05) & T-01 \\
RQ-1 & H1.3 & Q1.4 & $\kappa$ tutor--teacher vs.\ teacher--teacher; $\alpha$ (metrics module) & T-02 \\
RQ-1 & --- & Q1.5 & precision and recall by error type and letter & R-08 \\
RQ-2 & --- & Q2.1 & frequency of error kinds (SQL report 03) & T-03 \\
RQ-2 & H2.1 & Q2.2, Q2.4 & confusion matrices, all and confirmed (SQL reports 04, 10, 11) & R-09 \\
RQ-2 & --- & Q2.3 & difficult lines and minimal pairs (SQL reports 02, 08, 09) & T-04 \\
RQ-2 & H2.2 & --- & errors by articulatory attribute & R-10 \\
RQ-3 & H3.1 & Q3.1--Q3.2 & score change and recurrence after feedback & R-11 \\
RQ-3 & H3.2 & Q3.4 & delayed-test score & T-05 \\
RQ-4 & --- & Q6.1--Q6.3 & blocks by reason; interventions; iterations (journals) & R-13, T-08 \\
\bottomrule
\end{tabularx}
\end{table}

\section{Requirements traceability}

The project keeps a matrix that links each requirement to code and tests; an automated report regenerates the
current status and the fingerprint of the code version. At the time of writing 59 of 60 traced requirements
are covered by passing tests. The full matrix is part of the project repository and is reproduced in the
electronic appendix.

\chapter{Transcription for Russian speakers}\label{app:translit}

The tutor's hints use a simplified form of the scholarly Arabic--Russian transcription \cite{SRC-051}, so that
sounds absent from Russian remain distinct.

\begin{table}[ht]
\centering\small
\caption{Selected letters and their transcription in the tutor.}
\begin{tabular}{@{}cl@{\qquad}cl@{\qquad}cl@{}}
\toprule
Letter & Tr. & Letter & Tr. & Letter & Tr. \\
\midrule
\ar{ح} & \textrussian{х̣} & \ar{ص} & \textrussian{с̣} & \ar{ث} & \textrussian{с̱} \\
\ar{خ} & \textrussian{х} & \ar{ض} & \textrussian{д̣} & \ar{ذ} & \textrussian{з̱} \\
\ar{ع} & \textrussian{‘} & \ar{ط} & \textrussian{т̣} & \ar{ق} & \textrussian{к̣} \\
\ar{ء} & \textrussian{’} & \ar{ظ} & \textrussian{з̣} & \ar{غ} & \textrussian{г̣} \\
\bottomrule
\end{tabular}
\end{table}

A long vowel is written with two letters (\textrussian{аа, ии, уу}); syllables are ``consonant + vowel'', and
\emph{sukun} closes a syllable.

\chapter{Preregistration summary}\label{app:prereg}

The preregistration (version 0.2, to be approved and frozen before data collection) fixes:
\begin{itemize}
  \item seven confirmatory hypotheses (Table~\ref{tab:hyp}); everything else is exploratory;
  \item the unit of analysis (attempt) and clustering (learner);
  \item exclusions: stub attempts, unusable and off-task recordings (counted separately), attempts before the
    freeze; self-corrections in a separate block;
  \item participants: adults with consent to voice recording; first language and reading experience from the
    entry questionnaire;
  \item stopping rule: the planned end of the pilot; interim reports descriptive only;
  \item analysis: cluster bootstrap intervals, mixed models, Holm correction at $\alpha=0.05$;
  \item freezing: SHA-256 fingerprint and date recorded in the research log; any change requires a new version
    with a reason.
\end{itemize}
